\documentclass[11pt,a4paper]{article}

\usepackage[utf8]{inputenc}
\usepackage[T1]{fontenc}
\usepackage[spanish,es-noshorthands]{babel}
\usepackage{xcolor}
\usepackage[most]{tcolorbox}
\usepackage{listings}
\usepackage{fancyhdr}
\usepackage{lastpage}
\usepackage{enumitem}
\usepackage[hidelinks]{hyperref}

\setlength{\headheight}{30pt}
\pagestyle{fancy}
\fancyhf{}
\renewcommand{\headrulewidth}{0pt}
\fancyhead[R]{%
  \fontsize{8}{9.6}\selectfont
  Licenciatura en Ciencias de la Computación\\
  Sistemas Operativos\\[2pt]
  \rule{4.2cm}{0.5pt}
}
\fancyfoot[C]{\thepage\ / \pageref{LastPage}}

\newtcolorbox{infobox}{
  colback=blue!6!white, colframe=blue!45!white,
  arc=2mm, boxrule=0.6pt, left=8pt, right=8pt, top=6pt, bottom=6pt,
  title=INFO, fonttitle=\bfseries, coltitle=blue!45!black,
  colbacktitle=blue!6!white, boxsep=1pt, breakable
}

\newtcolorbox{extrabox}{
  colback=orange!8!white, colframe=orange!55!white,
  arc=2mm, boxrule=0.6pt, left=8pt, right=8pt, top=6pt, bottom=6pt,
  title=EXTRA, fonttitle=\bfseries, coltitle=orange!55!black,
  colbacktitle=orange!8!white, boxsep=1pt, breakable
}

\newtcolorbox{notebox}{
  colback=black!5!white, colframe=black!25!white,
  arc=2mm, boxrule=0.5pt, left=8pt, right=8pt, top=6pt, bottom=6pt,
  boxsep=1pt, breakable
}

\lstdefinestyle{codestyle}{
  backgroundcolor=\color{black!5},
  rulecolor=\color{black!25},
  frame=single,
  framerule=0.5pt,
  basicstyle=\ttfamily\small,
  breaklines=true,
  showstringspaces=false,
  tabsize=4,
  keywordstyle=\color{blue!60!black}\bfseries,
  commentstyle=\color{green!40!black},
  stringstyle=\color{orange!70!black},
  columns=fullflexible,
  upquote=true
}
\lstset{style=codestyle}
\setlist{nosep}

\title{Trabajo Práctico 4 -- Planificación}
\author{}
\date{}

\begin{document}
\maketitle
\thispagestyle{fancy}
\begin{center}
SISTEMAS OPERATIVOS
\end{center}
\newpage

\tableofcontents
\newpage

Este TP estudia cómo Linux decide qué tarea obtiene la CPU (Central
Processing Unit) y durante cuánto tiempo. Vas a observar prioridades,
políticas de planificación, cambios de contexto y desalojos usando comandos y
archivos del sistema. Los valores concretos dependen de la carga, la cantidad
de núcleos y la versión del núcleo: registrá tus observaciones en lugar de
copiar una salida de ejemplo.

\section{El planificador y las políticas de planificación}

El planificador (\emph{scheduler}) elige entre las tareas que están listas
para ejecutarse. Una tarea puede ceder la CPU voluntariamente al bloquearse,
por ejemplo al esperar entrada y salida (E/S, \emph{input/output}), o puede
ser desalojada cuando el núcleo necesita ejecutar otra tarea.

POSIX (\emph{Portable Operating System Interface}) define una interfaz para
consultar y modificar políticas de planificación. Linux implementa varias de
ellas; en este TP nos concentraremos en estas tres políticas y en la llamada
que permite ceder la CPU:

\begin{itemize}
  \item \texttt{SCHED\_OTHER}: es la política normal de tiempo compartido.
  Linux intenta repartir la CPU entre las tareas aptas y calcula una prioridad
  dinámica a partir de su comportamiento y de su valor de \texttt{nice}. No
  es conveniente modelarla como una planificación por turnos (\emph{Round
  Robin}) simple con un quantum (intervalo de tiempo, \emph{time slice}) fijo.
  \item \texttt{SCHED\_FIFO} (\emph{First In, First Out}): pertenece a las
  políticas de tiempo real y usa una prioridad estática. Una tarea continúa
  ejecutándose hasta bloquearse, ceder la CPU, terminar o ser desalojada por
  otra tarea de tiempo real con prioridad mayor. Las tareas con la misma
  prioridad no reciben un quantum automático.
  \item \texttt{SCHED\_RR} (\emph{Round Robin}): es similar a
  \texttt{SCHED\_FIFO}, pero las tareas de tiempo real con la misma prioridad
  se alternan mediante un intervalo de tiempo.
  \item \texttt{sched\_yield()}: no es una política adicional. Es una llamada
  que expresa que la tarea actual está dispuesta a ceder la CPU; su efecto
  depende de la política y de las demás tareas listas.
\end{itemize}

Las políticas de tiempo real pueden afectar a todo el sistema. En general,
un usuario común no puede asignárselas a sus procesos sin permisos especiales.
No las pruebes sobre procesos del sistema ni ejecutes una tarea de tiempo real
que nunca bloquee o termine en una máquina compartida.

\begin{lstlisting}[language=bash]
$ man 7 sched
$ ps -eo pid,ppid,ni,pri,rtprio,stat,psr,comm --sort=pid | head
$ chrt -p 1234
\end{lstlisting}

En la salida de \texttt{ps}, \texttt{NI} es el valor de \texttt{nice},
\texttt{PRI} es la prioridad que muestra \texttt{procps} y \texttt{RTPRIO}
es la prioridad de tiempo real cuando corresponde. \texttt{PSR} indica el
procesador lógico en el que se observó la tarea. La salida exacta y la
interpretación de algunas columnas dependen de la política y de la versión de
las herramientas.

\begin{extrabox}
Si tenés disponible el código fuente del núcleo de tu distribución, hacé esta
actividad; no hace falta leer todo el núcleo:

\begin{enumerate}
  \item Consultá la política del shell actual con
  \texttt{chrt -p \$\$} y anotá la salida.
  \item Ubicá la definición de \texttt{\_\_schedule()} con
  \texttt{grep -R -n \_\_schedule kernel/sched/}.
  \item Leé esa función e identificá dónde se obtiene la siguiente tarea (por
  ejemplo, mediante \texttt{pick\_next\_task}) y dónde se realiza el cambio de
  contexto (por ejemplo, mediante \texttt{context\_switch}).
  \item Entregá una tabla de dos filas y tres columnas: concepto descrito por
  \texttt{man 7 sched}, archivo y símbolo encontrados, y relación entre ambos.
  Las filas deben cubrir la selección de la siguiente tarea y el cambio de
  contexto. Si tu versión usa otros nombres o rutas, registrá la diferencia.
\end{enumerate}
\end{extrabox}

\section{Cambios de contexto}

Un cambio de contexto (\emph{context switch}) ocurre cuando el núcleo deja de
ejecutar una tarea y guarda su estado para poder continuar luego con otra. El
sistema de archivos virtual \texttt{/proc} (\emph{proc filesystem}) expone
parte de esta información como si fueran archivos.

El contador global de cambios de contexto aparece en la línea \texttt{ctxt}
de \texttt{/proc/stat}:

\begin{lstlisting}[language=bash]
$ grep '^ctxt ' /proc/stat
ctxt 123456789
\end{lstlisting}

Ese valor es acumulado desde el arranque y puede aumentar mientras observás
la salida. Repetí el comando varias veces: incluso una máquina que parece
inactiva ejecuta tareas de fondo, interrupciones y servicios.

Para estudiar un proceso particular, mirá su archivo \texttt{status}. Las
dos métricas siguientes separan las cesiones voluntarias de los desalojos
involuntarios:

\begin{lstlisting}[language=bash]
$ pid=$$
$ grep -E 'voluntary_ctxt_switches|nonvoluntary_ctxt_switches' \\
    /proc/$pid/status
voluntary_ctxt_switches:        4
nonvoluntary_ctxt_switches:     0
\end{lstlisting}

El contador de cesiones voluntarias aumenta cuando el proceso se bloquea o
cede la CPU de forma explícita. El contador de desalojos involuntarios aumenta
cuando el núcleo lo desaloja, por ejemplo para dar paso a otra tarea lista. No
esperes que uno de los dos sea siempre cero: un mismo programa puede hacer
ambas cosas.

\begin{notebox}
La variable especial \texttt{\$\$} contiene el PID (\emph{Process ID}) del
shell actual. En cambio, \texttt{/proc/self/status} se refiere al proceso que
abre el archivo. Si ejecutás \texttt{grep ... /proc/self/status}, normalmente
estarás observando al propio \texttt{grep}, que vive muy poco tiempo. Para
observar el shell u otro proceso, guardá su PID y construí
\texttt{/proc/\$pid/status} como en el ejemplo.
\end{notebox}

El desalojo no es sinónimo de hiperpaginación (\emph{thrashing}). La
hiperpaginación es un problema de memoria en el que el sistema pasa demasiado
tiempo moviendo páginas entre la memoria principal y el área de intercambio.
Muchos cambios de contexto pueden tener un costo de planificación, pero deben
analizarse junto con la carga de CPU y de memoria.

\begin{extrabox}
Si está instalado, compará el contador de \texttt{/proc} con
\texttt{pidstat -w 1} o \texttt{vmstat 1}. ¿Qué información es global y cuál
corresponde a un proceso? También podés explorar \texttt{/proc/interrupts},
sin modificar ningún archivo del sistema.
\end{extrabox}

\section{Observar procesos con \texttt{top}}

El comando \texttt{top} actualiza continuamente una lista de tareas y
permite cambiar su orden y enviarles señales. Probalo en una terminal:

\begin{lstlisting}[language=bash]
$ top
$ top -H
\end{lstlisting}

La opción \texttt{-H} muestra hilos (\emph{threads}) en lugar de agruparlos
por proceso. Dentro de \texttt{top}, las teclas más útiles para este TP son:

\begin{itemize}
  \item \texttt{P}: ordenar por porcentaje de CPU; \texttt{M}: ordenar por
  memoria;
  \item \texttt{1}: mostrar el uso de cada CPU lógica; \texttt{H}: alternar
  la vista de hilos;
  \item \texttt{k}: enviar una señal al PID elegido; \texttt{q}: salir.
\end{itemize}

Las columnas habituales son:

\begin{itemize}
  \item \texttt{PID}, \texttt{USER} y \texttt{COMMAND}: identificador,
  propietario y comando de la tarea;
  \item \texttt{PR} y \texttt{NI}: prioridad mostrada y valor de \texttt{nice}.
  En procesos normales, un valor de \texttt{NI} más alto representa menor
  preferencia relativa; los valores negativos requieren permisos adicionales;
  \item \texttt{VIRT}, \texttt{RES} y \texttt{SHR}: memoria virtual total,
  memoria residente y parte residente potencialmente compartida;
  \item \texttt{S} o \texttt{STAT}: estado. \texttt{R} es ejecutando o listo,
  \texttt{S} es sueño interrumpible, \texttt{D} es sueño ininterrumpible,
  \texttt{T} es detenido o trazado y \texttt{Z} es zombie;
  \item \texttt{\%CPU}, \texttt{\%MEM} y \texttt{TIME+}: uso reciente de
  CPU, proporción de memoria física y tiempo de CPU acumulado.
\end{itemize}

El orden de la lista puede cambiar aun cuando no escribas ningún comando.
El planificador actualiza prioridades dinámicas, aparecen y terminan tareas,
y en una máquina con varios procesadores distintas tareas pueden ejecutarse en
paralelo. Por eso una instantánea de \texttt{top} no alcanza para demostrar
una regla universal.

\section{Prioridad con \texttt{nice} y \texttt{renice}}

El valor de \texttt{nice} modifica la preferencia relativa de los procesos de
tiempo compartido. En Linux el rango usual va de $-20$ a $19$: un valor menor
representa más prioridad relativa y un valor mayor representa menos. El
comando \texttt{nice} crea el proceso con un incremento de niceness; el
comando \texttt{renice} modifica un proceso que ya existe.

\begin{lstlisting}[language=bash]
$ nice -n 10 comando
$ renice +5 -p 1234
\end{lstlisting}

En el primer ejemplo, \texttt{10} hace que el proceso sea menos preferido que
uno con \texttt{NI=0}. En el segundo, \texttt{+5} incrementa el valor de
niceness del PID indicado y también reduce su prioridad relativa. Bajar el
valor de niceness, por ejemplo a $-5$, normalmente requiere permisos
administrativos. No confundas \texttt{nice} con una medición de velocidad: si
la CPU está libre, dos procesos con valores distintos pueden usar casi toda
la CPU disponible.

\section{Experimento de planificación}

Vamos a comparar dos trabajos que consumen CPU. \texttt{yes} imprime una
secuencia interminable de letras; al redirigir su salida a \texttt{/dev/null}
la descartamos y dejamos visible principalmente el trabajo del procesador.
Abrí \texttt{top} en una terminal y usá otra para ejecutar lo siguiente:

\begin{lstlisting}[language=bash]
$ yes > /dev/null &
[1] 829
$ pid1=$!
$ nice -n 10 yes > /dev/null &
[2] 830
$ pid2=$!
\end{lstlisting}

Los números del ejemplo cambian en cada sistema. \texttt{\&} inicia el
comando en segundo plano (\emph{background}), \texttt{[1]} es el número de
trabajo del intérprete de comandos (\emph{shell}) y el otro número es su PID.
\texttt{\$!} guarda el PID del último proceso iniciado en segundo plano.
Observá ambos procesos en \texttt{top} durante varios
segundos y luego consultá sus métricas:

\begin{lstlisting}[language=bash]
$ ps -o pid,ppid,ni,pri,rtprio,stat,pcpu,comm -p "$pid1,$pid2"
$ grep -E 'voluntary_ctxt_switches|nonvoluntary_ctxt_switches' \\
    /proc/$pid1/status
$ grep -E 'voluntary_ctxt_switches|nonvoluntary_ctxt_switches' \\
    /proc/$pid2/status
\end{lstlisting}

Ahora cambiá la prioridad relativa del primer proceso, siempre hacia una
prioridad menor, y volvé a observar:

\begin{lstlisting}[language=bash]
$ renice +5 -p "$pid1"
$ ps -o pid,ni,pri,pcpu,stat,comm -p "$pid1,$pid2"
\end{lstlisting}

Registrá una tabla con los valores de \texttt{NI}, \texttt{PR}, \texttt{\%CPU}
y los dos contadores de contexto antes y después de \texttt{renice}. Contestá:

\begin{enumerate}
  \item ¿Qué cambió de manera inmediata y qué cambió sólo después de esperar?
  \item ¿El proceso con mayor niceness siempre recibió menos CPU? Repetí la
  prueba con poca carga y luego con dos o más procesos competidores.
  \item ¿Qué efecto tiene la cantidad de núcleos? Relacioná la respuesta con
  \texttt{PSR}, con la vista por CPU de \texttt{top} y con el hecho de que
  dos procesos pueden ejecutarse en paralelo.
  \item Compará los cambios de contexto voluntarios y no voluntarios. Proponé
  una explicación para cualquier diferencia, separando lo que observaste de
  lo que esperabas observar.
\end{enumerate}

Al terminar, finalizá los procesos usando una señal normal y esperá a que el
shell recoja sus estados:

\begin{lstlisting}[language=bash]
$ kill "$pid1" "$pid2"
$ wait "$pid1" "$pid2" 2>/dev/null
\end{lstlisting}

\begin{infobox}
\texttt{kill} envía una señal: no significa necesariamente ``matar''. La
señal predeterminada es \texttt{SIGTERM}, que permite a un programa terminar
ordenadamente. \texttt{Ctrl-C} envía normalmente \texttt{SIGINT} al trabajo en
foreground. Reservá \texttt{SIGKILL} (\texttt{kill -9}) para un proceso que no
responde y verificá siempre el PID antes de usarlo.
\end{infobox}

\section{Control de trabajos}

Como ya viste en el TP2, el intérprete de comandos distingue entre el primer
plano (\emph{foreground}), que ocupa la terminal, y el segundo plano
(\emph{background}), que permite volver al prompt. El número de trabajo sólo
tiene sentido dentro de ese shell; el PID identifica al proceso en el sistema.
Practicá el ciclo completo con un comando que tarde lo suficiente:

\begin{lstlisting}[language=bash]
$ sleep 100
\end{lstlisting}

Mientras está en foreground, apretá \texttt{Ctrl-Z}. Luego ejecutá:

\begin{lstlisting}[language=bash]
$ jobs
$ bg %1
$ jobs
$ fg %1
\end{lstlisting}

El número \texttt{1} es sólo un ejemplo: usá el que muestre tu shell. Con
\texttt{bg} el trabajo continúa en background y con \texttt{fg} vuelve al
foreground. Podés detenerlo otra vez con \texttt{Ctrl-Z} y terminarlo con
\texttt{kill \%1}. Relacioná cada acción con los estados que veas en
\texttt{jobs}, \texttt{ps} y \texttt{top}.

\begin{extrabox}
Explorá \texttt{taskset -p PID} para observar la afinidad de una tarea y
repetí el experimento usando una sola CPU, sólo si la máquina es tuya y sabés
que no vas a afectar a otros usuarios. Como extensión, investigá
\texttt{getrusage()} y la diferencia entre tiempo de usuario y tiempo de
sistema.
\end{extrabox}

\end{document}
